\chapter{实现与参数比较}

\label{ch:implementation}

\section{理论对象与实际代码}
\label{sec:implementation-table}

\begin{table}[H]
\centering
\caption{PMD、MDPO 与 BPO 的目标和实现差异}
\label{tab:method-comparison}
\begin{tabular}{P{0.16\textwidth}P{0.23\textwidth}P{0.23\textwidth}P{0.25\textwidth}}
\toprule
方法 & 理论更新 & 梯度/权重 & 主要偏差方差来源 \\
\midrule
PMD-Part & 回归 $u^+=\eta f-\log Z_f$ & 行为测度下 residual score，无外部 IS & 有限模型投影偏差；通常低于 IS 的 ratio 方差 \\
PMD-Mean & 中心化平方回归加 $\E_\beh[e^u]=1$ & Lambert-$W$ 投影梯度 & 均值和投影近似偏差；无乘积 IS \\
MDPO & 直接最大化 $\E_\pi[f]-\frac1\eta\KL(\pi\Vert\mu)$ & $\rho(f-\frac1\eta\log\rho)\nabla\log\pi$ & action IS 无偏但二阶矩随 $\E_\mu[\rho^2]$ 增长 \\
BPO & residual regression（本文件分析未线性化形式） & $\delta(g-\bar g)$，包含 score centering & $d_\beh$ 投影偏差与回归统计方差 \\
\bottomrule
\end{tabular}
\end{table}

\section{BPO 论文的训练预算}
\label{sec:bpo-implementation}

BPO 论文源码的主实验写明：每个 training step 先生成 256 个 prompt、每个 prompt 16 个 response，共 4096 条 response；再划分为 8 个、每个 512 条 response 的 mini-batch，执行 8 次 optimizer updates。主实验运行 400 个 training steps，即 3200 次 optimizer updates；附录消融运行 1000 个 training steps。论文没有按 residual norm 或 mirror-subproblem gap 停止的 inner loop，因此这些是固定更新预算，不是子问题收敛步数。

\section{OpenKimi/PMD 的训练预算}
\label{sec:openkimi-implementation}

OpenKimi 的 PMD loss 默认先对完整 response 的 token log-ratio 求和，再把 sequence log-ratio broadcast 到 response 内 token。\par Qwen2.5-7B math recipe 设置 `PPO\_EPOCHS=1`、`train\_batch\_size=512`、`actor.ppo\_mini\_batch\_size=32`，所以静态 batch 下每个 outer step 约有 $512/32=16$ 个 mini-batch 更新。代码没有 inner residual stopping 或 PMD gap 检验；`total\_training\_steps=3000` 是 outer steps 上限。

\section{经典 MDPO 的训练预算与 IS}
\label{sec:mdpo-implementation}

经典 MDPO on-policy 代码默认 `sgd\_steps=5`，MuJoCo README 的推荐示例使用 10。一个 rollout batch 固定后，对同一批数据重复执行该数量的 policy optimizer updates，再刷新 old/new anchor；没有根据 mirror residual 自适应停止。

off-policy 代码的 README 暴露 `sgd\_steps=1000`，但当前训练循环实际写为 `for grad\_step in range(1)`，原本的 `int(self.gradient\_steps)` 被注释。因此每个 train event 实际只执行一次 actor/critic update；`gradient\_steps` 主要影响 target/anchor policy 同步周期。实现使用 replay 状态测度和 action-level policy gradient，不能解释为完整的 prefix-IS 求解。

\section{比较时应报告的指标}
\label{sec:reporting}

为了比较四种方法，实验至少应固定同一 rollout batch，并记录：每个 inner update 的 mirror loss、log-ratio residual norm、$\KL(\cand_\theta\Vert\beh)$、梯度范数、ratio 的均值与二阶矩、中心化 residual 的条件方差，以及 prefix coverage。这样才能区分“做了更多 optimizer steps”和“真正接近 token-level mirror optimum”。
